\documentclass[a4paper]{article}
\usepackage{amsmath}
\usepackage{amssymb}
\usepackage{amsfonts}

\begin{document}

\normalsize

\noindent \large Auteur: Dina Haajou \\
\noindent \large Studierichting: Informatica \\
\noindent \large Oefeningenreeks 7A nr. 16.4 \\

\medskip

\normalsize

$f(x,y) = 
\left.
\begin{array}{rl}
\dfrac{x^2 y^8}{x^6 + y^{12}} & \text{voor } (x,y) \neq (0,0) \\
0 & \text{voor } (x,y) = (0,0)
\end{array}
\right.$ \\

Onderzoek naar continuïteit: \\

Laat $x = y^2$, dan wordt de limiet: \\

$= \lim\limits_{y \to 0} \dfrac{(y^2)^2 y^8}{(y^2)^3 + y^{12}} = \lim\limits_{y \to 0} \dfrac{y^{12}}{2y^{12}} = \dfrac{1}{2}$ \\

$\Rightarrow$ de limiet hangt af van het gekozen pad, dus $f$ is niet continu in $(0,0)$. \\

Onderzoek naar partiële afgeleiden: \\

$\dfrac{\partial f}{\partial x}(0,0) = \lim\limits_{\lambda \to 0} \dfrac{f(\lambda, 0) - f(0,0)}{\lambda} = \lim\limits_{\lambda \to 0} \dfrac{0}{\lambda^7} = 0$ \\

$\dfrac{\partial f}{\partial y}(0,0) = \lim\limits_{\lambda \to 0} \dfrac{f(0, \lambda) - f(0,0)}{\lambda} = \lim\limits_{\lambda \to 0} \dfrac{0}{\lambda^{13}} = 0$ \\

$\Rightarrow$ beide partiële afgeleiden bestaan en zijn nul in $(0,0)$. \\

Onderzoek naar afleidbaarheid: \\

$Df(0,h) = \lim\limits_{\lambda \to 0} \dfrac{f(\lambda h_1, \lambda h_2) - f(0,0)}{\lambda}$ \\

$= \lim\limits_{\lambda \to 0} \dfrac{\lambda^{10} h_1^2 h_2^8}{\lambda^7 (h_1^6 + \lambda^6 h_2^{12})}$ \\

$= \lim\limits_{\lambda \to 0} \dfrac{\lambda^3 h_1^2 h_2^8}{h_1^6 + \lambda^6 h_2^{12}} = 0$ \\

$\Rightarrow$ alle richtingsafgeleiden bestaan en zijn nul, dus $f$ is afleidbaar in $(0,0)$. \\

Onderzoek naar differentieerbaarheid: \\

Omdat de functie niet continu is in $(0,0)$ \\

volgt meteen dat $f$ niet differentieerbaar is in dat punt. \\

\begin{thebibliography}{999}
% \bibitem{}
\end{thebibliography}

\end{document}
